
@@M0@@
D'après la
proposition 6.11, l'homomorphisme : @@M1@@ est un isomorphisme si et seulement si l'homomorphisme : @@M2@@ en est un. Comme il est visiblement surjectif, ce dernier l'est aussi, et il suffit donc de montrer que @@M3@@ est injectif. Par fonctorialité de @@M4@@, on obtient le diagramme :
@@M5@@
comme @@M6@@ est à degrés bornés, la flèche verticale de droite est un isomorphisme (6.12), et en utilisant 6.11, on obtient donc un morphisme : @@M7@@, le composé des deux étant l'identité d'après la propriété universelle de @@M8@@.

\section*{APPENDICE : Les opérations @@M9@@ d'Adams}

Nous suivons ici une présentation due à J.P. Serre.

7.1. Soit @@M10@@ un pré-@@M11@@-anneau. Pour tout @@M12@@, on désigne par @@M13@@ la dérivée par rapport à @@M14@@ de @@M15@@. On définit alors des éléments de @@M16@@, notés @@M17@@, en posant :
@@M18@@

Les applications @@M19@@ sont additives~; en effet, on a :
@@M20@@
